\documentclass[12pt,a4paper]{article}
\usepackage[margin=0.75in]{geometry}
\usepackage{xcolor}
\usepackage{hyperref}
\usepackage{amsmath}
\usepackage{amssymb}
\usepackage{fancyhdr}

\pagestyle{fancy}
\fancyhf{}
\rhead{\thepage}
\lhead{Semantic Video Understanding}

\title{\textbf{Semantic Video Understanding and Object-Level Event Detection} \\
\Large{A Hybrid Approach Integrating YOLOv8 with CLIP Embeddings}}
\author{
    Prabal Minotra (2022357) \\
    Sarthak Bhudolia (2023491) \\
    Altamash Hasnain (2023073) \\
    Bhuvan Raju (2023171) \\
    \\
    \textit{Indraprastha Institute of Information Technology, Delhi}
}
\date{\today}

\begin{document}

\maketitle

\begin{abstract}

This report presents a semantic video understanding system that combines YOLOv8 object detection with CLIP-based semantic embeddings for robust video analysis. The system implements a hybrid two-tier detection strategy: a fast primary path using YOLO followed by CLIP verification, and a fallback mechanism employing sliding-window CLIP scans for detecting missed objects. We analyze the architectural design, algorithmic foundations, and practical applications. The approach achieves real-time performance on modern GPUs while maintaining semantic understanding capabilities, enabling natural language queries for video search and compliance monitoring in safety-critical applications.

\end{abstract}

\textbf{Keywords:} Object Detection, CLIP Embeddings, Video Analysis, Semantic Understanding, YOLOv8

% ============================================================================
\section{Introduction}
% ============================================================================

Video analysis at scale requires balancing two competing objectives: real-time performance and semantic understanding. Traditional object detection approaches (e.g., YOLO) provide speed but lack semantic context, while vision-language models (e.g., CLIP) offer semantic capabilities at higher computational cost. This work investigates a hybrid system that leverages both paradigms to enable efficient, interpretable video analysis supporting natural language queries.

The key contribution is a two-tier detection architecture that maintains computational efficiency while improving robustness. The primary tier uses YOLOv8 for rapid object detection followed by CLIP-based semantic verification. The secondary tier implements a multi-scale sliding-window CLIP scan as a fallback mechanism when primary detection fails. This design addresses three critical challenges: achieving interactive frame rates, supporting semantic understanding beyond predefined classes, and robustly detecting objects in challenging scenarios.

% ============================================================================
\section{System Architecture and Methodology}
% ============================================================================

\subsection{Hybrid Detection Strategy}

The system employs a decision-based detection pipeline operating in two modes:

\textbf{Tier 1 (Fast Path):} For each video frame, YOLOv8 generates object candidates with bounding boxes and class labels. For each detection, a 25\% contextual padding is applied to create image crops suitable for CLIP processing. The CLIP model encodes both the natural language query and each image crop to normalized 512-dimensional embeddings. Cosine similarity between query and image embeddings determines relevance. A label-based semantic boost is applied: if the YOLO label overlaps with query terms, the similarity score increases by 0.08 (direct match) or 0.05 (partial match). Candidates exceeding the similarity threshold are returned as matches.

\textbf{Tier 2 (Fallback Path):} When Tier 1 produces no matches, a multi-scale sliding-window scan covers the entire frame. The scan operates at three scales (0.15, 0.10, 0.08 of frame dimensions) with a systematic 6×8 grid pattern. At each window position, CLIP embeddings are computed and compared against the query. A novel zero-shot filtering mechanism uses negative prompts (e.g., "background," "wall," "shadow") to suppress false positives. Candidates are accepted only if the query probability exceeds 40\% in a softmax distribution over all prompts. Non-maximum suppression with IoU threshold 0.4 removes duplicate detections.

The decision criterion is:

\begin{equation}
\text{result} = \begin{cases}
\text{Tier 1 matches} & \text{if } \#\text{matches} > 0 \\
\text{Tier 2 scan} & \text{otherwise}
\end{cases}
\end{equation}

\subsection{Temporal Processing}

Detected object timestamps are aggregated into coherent temporal segments. The aggregation algorithm groups detections based on temporal proximity: consecutive detections within a 2.0-second gap are merged into a single segment, while larger gaps initiate new segments. This enables temporal event summarization and efficient compliance checking.

Color information is extracted from each detection by analyzing the HSV color space of object crops. A median-based approach focusing on the central 50\% of crops provides robustness against background contamination. Objects are classified into twelve color categories (red, orange, yellow, green, cyan, blue, purple, pink, gray, black, white) based on HSV thresholds. This adds semantic context to detections and enables color-based search queries.

Compliance checking operates on the detection history. Rule-based logic identifies violations: for instance, a person detected without an accompanying helmet triggers a safety violation. This design allows straightforward extension with additional rules for specific domains.

% ============================================================================
\section{Results and Analysis}
% ============================================================================

\subsection{Experimental Findings}

CLIP model validation demonstrates stable semantic embeddings. Testing on representative objects (yellow\_box, door\_shadow, actual\_broom) yielded average raw similarities of 0.231 with low variance (0.008), indicating consistent behavior. Notably, background elements (e.g., shadows) scored lower (0.220) than actual objects (0.240), confirming discriminative capability.

Performance benchmarks on RTX 3090 GPU show YOLOv8 nano processes frames at 5 milliseconds per frame ($\approx$200 FPS), CLIP single-image encoding at 3 milliseconds, and query matching at 15 milliseconds for 5 detections. The full Tier 1 pipeline achieves approximately 30 FPS. Tier 2 sliding-window scanning costs 500 milliseconds per frame, making it suitable for offline or batch processing scenarios.

Computational analysis reveals Tier 1 operates in $O(1)$ time per frame (bounded by model size), while Tier 2 complexity scales as $O(s \cdot g^2)$ where $s=3$ scales and $g=6-8$ grid dimension, resulting in roughly 108 CLIP evaluations per frame. This trade-off between speed and robustness is configurable via parameter thresholds.

\subsection{Design Rationale}

The label-based semantic boost improves matching accuracy by leveraging complementary information: YOLO provides precise object category information, while CLIP captures fine-grained visual and semantic properties. For instance, distinguishing between different car models becomes possible by combining YOLO's "car" classification with CLIP's ability to match specific colors and designs.

The zero-shot filtering mechanism addresses a fundamental challenge in sliding-window scanning: distinguishing true objects from background regions. By using negative prompts representing common non-target patterns (walls, floors, shadows), the system suppresses approximately 95\% of false positives while maintaining reasonable recall. The 0.40 probability threshold balances sensitivity against specificity.

The temporal aggregation with 2.0-second gap threshold handles typical detection variability. Brief interruptions in detection (due to occlusion or transient fails) are grouped into continuous events, while longer gaps correctly identify distinct object instances. This parameter is application-dependent and tunable.

\subsection{Practical Applications}

The system enables three primary use cases: (1) semantic video search via natural language queries (e.g., "find all red cars"), (2) safety compliance monitoring by detecting violations (e.g., unhelmed workers in construction sites), and (3) temporal event analysis by identifying when and for how long objects appear. The ability to support arbitrary queries without retraining makes CLIP-based approaches particularly valuable for dynamic scenarios where object types are not known a priori.

% ============================================================================
\section{Limitations and Conclusions}
% ============================================================================

Current limitations include: (1) lack of multi-object tracking across frames, (2) reliance on hand-crafted compliance rules, (3) computational expense of Tier 2 for real-time applications, and (4) performance degradation under extreme lighting or occlusion. Future work should explore learned compliance rules via neural networks, optimized spatial-temporal scanning strategies, and multi-object trajectory matching using Kalman filtering or Hungarian algorithms.

This work demonstrates that hybrid approaches effectively combine the strengths of classical detection and semantic embeddings. The two-tier architecture provides a practical framework for semantic video understanding, achieving computational efficiency on Tier 1 while maintaining robustness through Tier 2. The system successfully addresses key challenges in video analysis: speed, semantic understanding, and robustness. Applications in safety monitoring and intelligent surveillance demonstrate practical value. The modular design enables straightforward extension with additional modules for specific domains or requirements.

\end{document}